\section{Resident count predicates and delayed device feedback}
\label{sec:r15-resident-gpu}
The resident profile keeps point seeds, a derived hierarchy and a query batch
in device memory across selected launches. Its output is a count and a
resolution status. This is a different interface from returning every point
identity; comparisons must use the same requested output. A separate
working-set experiment evaluates delayed exact word feedback. Its timing
does not by itself measure a spherical-query improvement.

\subsection{The exact distance target and a bounded binary64 body}
For the original finite binary64 triples $p,q$ accepted by the pinned S2
unit-length predicate, define
\begin{equation}
 P=p/\|p\|,\qquad Q=q/\|q\|,\qquad
 D=\|P-Q\|^2,\qquad 0\le D\le4.
 \label{eq:r15-resident-target}
\end{equation}
The supplied binary64 radius $R^2\in[0,4]$ denotes its exact dyadic value;
membership is $D\le R^2$. The stored triples are not silently assumed to
have exact norm one. This normalized-direction meaning is the one declared
by the official
\href{https://github.com/google/s2geometry/blob/079611b654ad89afd9c3c3a1796d64bdd6a6b340/src/s2/s2predicates.h}
{S2 distance predicate}.

Let $\varepsilon=2^{-52}$ and assume strict binary64 round-to-nearest,
gradual underflow and the displayed non-fused operation sequence. S2's
computed norm tolerance is $5\varepsilon$. Its three-square evaluation error
gives $|\|p\|^2-1|<9\varepsilon$, hence
$|\|p\|-1|<5\varepsilon$ and the conservative coordinate bound
$|p_i-P_i|<8\varepsilon$. Apply the same bound to $q$.

The direct device body is
\begin{equation}
 d_i=\operatorname{RN}(p_i-q_i),\quad
 s_i=\operatorname{RN}(d_i^2),\quad
 \widehat D=\operatorname{RN}
       \bigl(\operatorname{RN}(s_1+s_2)+s_3\bigr).
 \label{eq:r15-resident-distance}
\end{equation}
Writing $\Delta_i=P_i-Q_i$, raw coordinate error plus subtraction rounding
gives $|d_i-\Delta_i|<18\varepsilon$. Since $|\Delta_i|\le2$,
the square error before rounding is below
$18\varepsilon(4+18\varepsilon)<73\varepsilon$; product rounding adds less
than $3\varepsilon$. Each of the two additions adds less than
$8\varepsilon$. Therefore
\begin{equation}
 |\widehat D-D|<3(76\varepsilon)+2(8\varepsilon)
             =244\varepsilon<256\varepsilon.
 \label{eq:r15-resident-error}
\end{equation}
The detailed proof includes the elementary absolute underflow allowance
$2^{-1075}$ and the norm-admission bound in
\code{formal/RESIDENT\_PREDICATE\_PROOF.md}. No overflow is possible in this
bounded coordinate body. This is an error certificate in dimensionless
squared-chord units, not a physical position tolerance.

\subsection{A wide interval produces definite or unresolved results}
Use the exactly representable margin $m=2^{-40}=4096\varepsilon$ and
\begin{equation}
 t_-=\operatorname{RN}(R^2-m),\qquad
 t_+=\operatorname{RN}(R^2+m).
 \label{eq:r15-resident-thresholds}
\end{equation}
Each threshold's rounding error is below $3\varepsilon$, including zero,
subnormal radii and $R^2=4$. Thus
\begin{equation}
 \begin{array}{lll}
 \widehat D\le t_-&\Longrightarrow D<R^2&\text{definitely inside},\\
 \widehat D>t_+&\Longrightarrow D>R^2&\text{definitely outside},\\
 \text{otherwise}&&\text{unresolved by this interval}.
 \end{array}
 \label{eq:r15-resident-classification}
\end{equation}
For example, the first implication follows from
$D<\widehat D+256\varepsilon
 <R^2-(4096-3-256)\varepsilon$ after the threshold comparison.
The second follows from the corresponding lower bound. An exact boundary
cannot be assigned to outside by either proved rule.

The resident implementation uses componentwise equality of the original
triples as an exact
zero-distance certificate, including signed zeros. A rounded distance equal
to zero is not: distinct subnormal separations can square to zero. Distinct
proportional triples may also denote the same direction and require an
appropriate exact certificate or unresolved result. The endpoint $R^2=4$
separately accepts every admitted direction from
equation~\eqref{eq:r15-resident-target}.

\subsection{Complete counts and a visible resolution path}
The conservative binary32 hierarchy filter can reject only proved outside
subtrees. Surviving raw binary64 point words reach the predicate above. If
$I$ points are definitely inside and $U$ remain unresolved, then
\begin{equation}
 I\le\operatorname{Count}(q,R^2)\le I+U.
 \label{eq:r15-resident-count}
\end{equation}
Only $U=0$ establishes a complete resident count. Otherwise the query status
remains unresolved until a selected exact host query supplies the whole
answer. Partial counts are not final answers. The count profile does not
truncate an ID list or inherit a candidate-capacity prefix as its result.

Seeds, query words and bounds share the selected epoch. Derived outputs use
separate storage, and their consumers wait for completed producer launches.
An immutable point hierarchy can be reused; changed geometry needs a new
valid bound and binding. Texture instructions do not guarantee a cache hit
or that every resident VRAM word is resident in the texture cache.

\subsection{The implemented count-to-word feedback binding}
The \code{ResidentFeedback} profile assigns immutable query alternatives
$\mathcal Q_i^0,\mathcal Q_i^1$ to lane $i$. A count kernel reads the old
64-bit state word $q_i^n$ and evaluates
\begin{equation}
 (I_i^n,U_i^n)=\operatorname{ResidentCount}
       \bigl(\mathcal Q_i^{\,q_i^n\mathbin{\&}1}\bigr).
 \label{eq:r15-resident-selection}
\end{equation}
A following commit kernel consumes that summary. If $U_i^n>0$, its word,
parity, hinge count and last committed epoch hold; the status is unresolved.
This feedback profile does not silently run a host resolver. For a resolved
summary, the drive is $d_i^n=[I_i^n>0]$. A new resolved epoch uses the explicit
pulse-on-first rule
\begin{equation}
 h_i^n=[\text{uninitialized}]
        \mathbin{\lor}[d_i^n\ne d_i^{\rm last}].
 \label{eq:r15-resident-pulse}
\end{equation}
Replaying the last committed epoch holds; an older epoch is rejected.
Only resolved, ordered epochs update the remembered drive and epoch.
Thus the first sample is a declared initialization event, not an inferred
geometric crossing.

Let $T_x,T_J,T_K$ be the supplied bitwise four-entry Boolean tables, let $M$
be the valid-word mask, and write the two ordered ASA/NA stages as
\begin{equation}
 \mathcal A_j(v)=
 \begin{cases}
 0,&(v\mathbin{\&}A_j\mathbin{\&}N_j)\mathbin{\&}B_j\ne0,\\
 v\mathbin{\&}A_j\mathbin{\&}N_j,&\text{otherwise}.
 \end{cases}
 \label{eq:r15-resident-mask}
\end{equation}
For an accepted pulse, all three tables refer to the same old word:
\begin{align}
 y&=\mathcal A_1\bigl(\mathcal A_0(T_x(q,d)\mathbin{\&}M)\bigr),\\
 J&=T_J(q,y)\mathbin{\&}M,\qquad K=T_K(q,y)\mathbin{\&}M,\\
 q^+&=((J\mathbin{\&}\neg q)\mathbin{\lor}
                      (\neg K\mathbin{\&}q))\mathbin{\&}M,
 \qquad p^+=p\mathbin{\oplus}1.
 \label{eq:r15-resident-jk}
\end{align}
Otherwise the logical state holds. Exhausted hinge-count storage rejects the
commit. This profile has a fixed chart, so its orientation stays zero.

The source and query table outlive the feedback object. Each external phase
acquires a fresh invalidated device view, and calls sharing that source are
serialized. The count kernel reads state; the distinct commit kernel writes
it. Their common stream orders these operations before the next count kernel,
so a second state buffer is unnecessary for this two-kernel binding. No
per-epoch host data readback is required. This is delayed feedback through a
geometric predicate, with unresolved results still visible.

\subsection{A separate exact integer working-set profile}
The working-set executable binds a smaller independent operator profile.
Each immutable 16-byte record holds integer $a,b\in[-1023,1023]$, an opcode
$o\in\{0,1,2,3\}$ and a 32-bit salt $w$. Its old state is a 32-bit word $q$
and a parity bit. Define
\begin{equation}
 s=(q\mathbin{\&}31)-16,\qquad
 c=(q\mathbin{\&}3)-1,\qquad
 d=((q\gg2)\mathbin{\&}3)-1,
\end{equation}
where subtraction occurs in a signed integer domain. With
$\phi=(1+\sqrt5)/2$, the selected exact coefficient pair is
\begin{equation}
 (a',b')=
 \begin{cases}
 (a+s,b),&o=0,\\
 (b,a+b),&o=1,\\
 (ac+bd,\ ad+bc+bd),&o=2,\\
 (a-s,b),&o=3.
 \end{cases}
 \label{eq:r15-working-phi}
\end{equation}
The third row is literal multiplication by $c+d\phi$, using
$\phi^2=\phi+1$. No floating-point approximation to $\phi$ is evaluated.
The bounds $|a'|\le4092$, $|b'|\le6138$ and
$|2a'+b'|\le14322$ give
\begin{equation}
 (2a'+b')^2\le205119684,\qquad
 5b'^2\le188375220<2^{31}-1.
 \label{eq:r15-working-sign32}
\end{equation}
Both squares and every preceding coefficient temporary fit signed 32-bit
arithmetic. The optimized device body uses that width without narrowing loss;
the independent CPU recurrence retains 64-bit intermediates.
Same-sign terms are handled directly; the square comparison is used only
for opposing signs in $2(a'+b'\phi)=2a'+b'+b'\sqrt5$.
This proof is specific to the small procedural coefficient domain above.
It does not narrow the separate native expression profile whose coefficients
may approach $2^{30}-1$.

For $\sigma=\operatorname{sign}(a'+b'\phi)=0$ the state holds. Otherwise
the drive word is $w$ for $\sigma<0$ and $\neg w$ for $\sigma>0$;
$x=q\mathbin{\oplus}\text{drive}$ passes through the two fixed mask stages
in equation~\eqref{eq:r15-resident-mask}. The selected $J=y,K=\neg y$
specialization gives $q^+=y$ and toggles parity. Every launch epoch supplies
one fresh internal event opportunity per record; repeated nonzero signs do
not imply repeated physical side crossings. Old $q$ changes the next
operands and word body, making this an explicit delayed self-reference.

\subsection{Bank coverage and coherence}
The immutable seed records are expanded from integer bit planes through raw
integer texture reads. In this executable, only immutable records use the
texture/global comparison; the two state buffers use global memory in both
variants. Every record reads its epoch's old buffer and writes the distinct
next buffer. All chunks finish before buffer roles change. The packed scratch
producer and its texture consumer are also separate completed launches.
These boundaries follow the ordering and texture-coherence requirements in
the official
\href{https://docs.nvidia.com/cuda/archive/12.8.1/cuda-c-programming-guide/index.html#streams}
{CUDA stream} and
\href{https://docs.nvidia.com/cuda/archive/12.8.1/cuda-c-programming-guide/index.html#read-write-coherency}
{texture-coherence} documentation; they do not assume coherent simultaneous
texture reads and global writes to the same address.

For $N$ records, streaming visits ordinal $k$ directly. The other pattern uses
\begin{equation}
 i(k)=(\alpha k+\beta)\bmod N,\qquad
 \gcd(\alpha,N)=1,\qquad 0\le k<N.
 \label{eq:r15-working-permutation}
\end{equation}
Checked integer products avoid overflow, and coprimality makes this a
bijection. Thus each full epoch visits every record exactly once, including
every bank, rather than probing a small subset of a large allocation.
The declared working set contains $16N$ seed bytes and $16N$ bytes in the two
persistent state buffers. Its size is admitted against available VRAM with
a reserve; only completed runs establish an exercised size. This accounting
does not mean that all VRAM allocations fit in hardware caches.

\subsection{Exact address and digest implementation}
The optimized address body preserves equation~\eqref{eq:r15-working-permutation}.
Streaming uses $i=k$ directly. For a power-of-two record count, the remainder
is the low-bit mask. For every other admitted $N\ge2048$, put $W=2^{64}$
and, for the checked numerator $0\le x=\alpha k+\beta<W$, use
\begin{equation}
 m=\left\lfloor\frac WN\right\rfloor,\qquad
 q_0=\left\lfloor\frac{xm}{W}\right\rfloor,\qquad
 r=x-q_0N,\qquad
 i=\begin{cases}r-N,&r\ge N,\\r,&r<N.\end{cases}
 \label{eq:r15-working-reciprocal}
\end{equation}
The device obtains $q_0$ from the high half of an exact unsigned 64-bit
product. Since
$0\le x/N-xm/W\le x/W<1$, $q_0$ is either
$\lfloor x/N\rfloor$ or one less. Hence $0\le r<2N$ and one subtraction
produces the exact remainder. The implementation computes $m$ on the host
from $W-1$, avoiding a stored integer equal to $W$; the existing checked
affine product and addition remain required. This replaces division, not the
permutation or its coverage contract.

Each full bank has a verified capacity $C=2^b$ within the device texture
element limit and the selected cap. Its index and offset are exactly
$i\gg b$ and $i\mathbin{\&}(C-1)$. The final bank retains its actual valid
count. Both immutable-record texture and global-load variants use the same
address mapping, operands, recurrence and old/new state allocation rules.

The 128-thread digest now reduces four warp summaries, then combines those
summaries in the first warp after one block barrier. All lanes participate;
invalid ordinals contribute zero XOR, zero sum and zero count. XOR and
unsigned addition modulo $2^{64}$ are associative, and the block count is at
most 128, so regrouping changes no digest bits. The implementation contract
and independent-oracle limits are recorded in
\href{formal/WORD_CACHE_BENCHMARK.md}{\code{formal/WORD\_CACHE\_BENCHMARK.md}}.
These exact transformations provide no speed claim without the corresponding
compiled checks and matched measurements.

\subsection{Validation and measurement scope}
Paired texture/global trials reset both state buffers to the same procedural
snapshot outside timing. Per-epoch validation combines full-record XOR, sum
and count digests with sampled independent CPU state replay. For
$N\le65536$, an additional full CPU digest is computed. Digests can collide;
they are evidence distinct from exhaustive bitwise state comparison. The
timed kernel includes arithmetic, masks, JK and digest work. The reported
$32N$ logical bytes per epoch count seed read, old-state read and next-state
write; their effective rate is not a hardware DRAM-bandwidth measurement.

The count predicate and this word recurrence have different outputs and
costs. Their measurements keep resident kernel time, transfers, host
resolution, matched global loads and validation overhead separate. Neither a
large working set nor their shared texture path establishes a speed advantage
for the other's task, a cache-hit guarantee or physical navigation accuracy.

\IfFileExists{resident_results.tex}{% Generated only from complete, validated reports by tools/summarize_resident.py.

\section{Measured resident counts and word working sets}

\label{sec:resident-results}

All 27 resident count cases were measured; every complete answer matched.

The texture path with complete host-visible counts beat the best measured S2 configuration in 20 of 27 cases and the best measured CPU implementation in 13 of 27 cases.

These results concern inclusive spherical \emph{counts} over the supplied exact IEEE seeds. They are not ID-return, general-S2, physical-position-accuracy or photonics measurements.

The word-cache results below are the preserved baseline in \path{review/word_cache_full.json}; later optimizations require separate measurement records.

\subsection{Every resident count case}

All entries are median milliseconds; CPU cells show the chosen worker count in parentheses. T/G mean integer texture/global loads. Device time is per resident launch, with five launches per sample. Complete time includes launch, count readback and any host exact fallback. The last column is unresolved points / texture fallback queries / global fallback queries.

\begingroup
\scriptsize
\setlength{\tabcolsep}{2.5pt}
\renewcommand{\arraystretch}{1.16}
\begin{longtable}{lrrrrrrrrr}
\toprule
Family & $N$ & $Q$ & S2 best & BVH best & T dev. & T comp. & G dev. & G comp. & U/Ft/Fg \\
\midrule\endfirsthead
\toprule
Family & $N$ & $Q$ & S2 best & BVH best & T dev. & T comp. & G dev. & G comp. & U/Ft/Fg \\
\midrule\endhead
\bottomrule\endfoot
uniform & 16384 & 256 & 0.1165 (20) & 0.0476 (20) & 0.2154 & 0.2685 & 0.2031 & 0.2571 & 0/0/0 \\
uniform & 16384 & 4096 & 1.038 (20) & 0.3499 (20) & 0.2341 & 0.3236 & 0.2237 & 0.3168 & 0/0/0 \\
uniform & 16384 & 65536 & 11.42 (20) & 3.458 (20) & 1.304 & 1.776 & 1.276 & 1.826 & 0/0/0 \\
uniform & 262144 & 256 & 0.1888 (20) & 0.0799 (20) & 0.2536 & 0.3096 & 0.2447 & 0.2975 & 0/0/0 \\
uniform & 262144 & 4096 & 1.52 (20) & 0.4807 (20) & 0.2771 & 0.3811 & 0.2642 & 0.3538 & 0/0/0 \\
uniform & 262144 & 65536 & 19.96 (20) & 6.063 (20) & 1.587 & 2.15 & 1.552 & 2.007 & 0/0/0 \\
uniform & 1048576 & 256 & 0.2667 (20) & 0.1105 (20) & 0.3203 & 0.3787 & 0.3051 & 0.3612 & 0/0/0 \\
uniform & 1048576 & 4096 & 2.365 (20) & 0.628 (20) & 0.3155 & 0.4242 & 0.3005 & 0.3899 & 0/0/0 \\
uniform & 1048576 & 65536 & 28.74 (20) & 6.631 (20) & 2.793 & 3.913 & 2.754 & 3.92 & 0/0/0 \\
clustered & 16384 & 256 & 1.423 (20) & 0.234 (20) & 1.529 & 1.629 & 1.428 & 1.531 & 0/0/0 \\
clustered & 16384 & 4096 & 13.8 (20) & 2.232 (20) & 1.573 & 1.847 & 1.472 & 1.73 & 0/0/0 \\
clustered & 16384 & 65536 & 213.4 (20) & 26.94 (20) & 12.57 & 13.1 & 12.36 & 12.85 & 0/0/0 \\
clustered & 262144 & 256 & 3.168 (20) & 0.4402 (20) & 8.505 & 8.633 & 8.06 & 8.193 & 0/0/0 \\
clustered & 262144 & 4096 & 35.02 (20) & 4.232 (20) & 8.864 & 8.964 & 8.366 & 8.503 & 0/0/0 \\
clustered & 262144 & 65536 & 548.8 (20) & 66.61 (20) & 64.18 & 64.95 & 63.23 & 63.95 & 2/2/2 \\
clustered & 1048576 & 256 & 4.473 (20) & 0.5397 (20) & 18.65 & 18.79 & 17.76 & 17.9 & 0/0/0 \\
clustered & 1048576 & 4096 & 52.52 (20) & 6.573 (20) & 18.65 & 18.79 & 17.77 & 17.9 & 0/0/0 \\
clustered & 1048576 & 65536 & 744 (20) & 76.54 (20) & 132.3 & 133.3 & 130.4 & 131.5 & 7/7/7 \\
great\_circle & 16384 & 256 & 0.563 (20) & 0.0843 (4) & 0.2874 & 0.3448 & 0.272 & 0.3235 & 0/0/0 \\
great\_circle & 16384 & 4096 & 3.613 (20) & 0.502 (20) & 0.3082 & 0.4366 & 0.2926 & 0.4139 & 0/0/0 \\
great\_circle & 16384 & 65536 & 42.77 (20) & 5.61 (20) & 2.04 & 2.505 & 2.002 & 2.438 & 0/0/0 \\
great\_circle & 262144 & 256 & 1.495 (20) & 0.1508 (20) & 1.138 & 1.25 & 1.082 & 1.197 & 0/0/0 \\
great\_circle & 262144 & 4096 & 10.49 (20) & 1.129 (20) & 1.142 & 1.41 & 1.09 & 1.341 & 0/0/0 \\
great\_circle & 262144 & 65536 & 175.9 (20) & 20.45 (20) & 7.901 & 8.489 & 7.776 & 8.395 & 0/0/0 \\
great\_circle & 1048576 & 256 & 2.122 (20) & 0.2988 (20) & 2.943 & 3.098 & 2.79 & 2.975 & 0/0/0 \\
great\_circle & 1048576 & 4096 & 25.14 (20) & 3.806 (20) & 3.935 & 4.06 & 3.779 & 3.9 & 0/0/0 \\
great\_circle & 1048576 & 65536 & 380.9 (20) & 44.99 (20) & 21.58 & 22.19 & 21.23 & 21.78 & 0/0/0 \\
\end{longtable}
\endgroup

There were 9 unresolved point decisions across 2 cases; 9 queries per access mode required complete host fallback. Texture beat global loads in 0 of 27 device-time comparisons and 2 of 27 complete-time comparisons. The GPU-over-CPU gain therefore does not establish a texture-cache advantage.

The first GPU construction cost 2547.1377 ms, including 2546.8671 ms reported source preparation/upload. This cold setup is excluded from repeated-query timings; later setup costs remain in the complete report.

Device-only counts are complete only where the unresolved count is zero. Every 1/4/20-worker configuration and raw trial is retained in \path{review/resident_summary.json}; best means the lowest measured median, not an unreported preferred configuration.

S2 uses its public closest-point API with reused candidate vectors, followed by exact refinement/counting. The BVH count avoids ID allocation. Setup is excluded. GPU device-mode order alternates; CPU/configuration and host-complete mode blocks have fixed order. Five warm trials provide no confidence interval or universal speed guarantee.

\subsection{Word/state coverage through the memory hierarchy}

NVIDIA GeForce RTX 5070 Ti Laptop GPU reports 36 MiB of L2 (37748736 bytes). The largest working set was 9.744 GiB, 277.2 times L2 and 81.61\% of reported device memory. The configured VRAM reserve remained available.

This is a separate bounded integer phi-expression and ASA/NA+JK word/state workload, \emph{not an S2 comparison}. Working bytes include immutable records plus both persistent state buffers. Every record is visited in every epoch; zero-sign holds may leave a value unchanged.

\begingroup
\scriptsize
\setlength{\tabcolsep}{2.5pt}
\renewcommand{\arraystretch}{1.16}
\begin{longtable}{rrrrrl}
\toprule
Working MiB & L2 multiple & Records/epoch & Allocated MiB & Sampled & CPU oracle \\
\midrule\endfirsthead
\toprule
Working MiB & L2 multiple & Records/epoch & Allocated MiB & Sampled & CPU oracle \\
\midrule\endhead
\bottomrule\endfoot
0.0625 & 0.001736 & 2048 & 0.09888 & 124 & Full replay \\
0.25 & 0.006944 & 8192 & 0.3815 & 130 & Full replay \\
1 & 0.02778 & 32768 & 1.511 & 131 & Full replay \\
4 & 0.1111 & 131072 & 6.028 & 131 & Sample replay \\
16 & 0.4444 & 524288 & 24.05 & 131 & Sample replay \\
18 & 0.5 & 589824 & 26.05 & 131 & Sample replay \\
36 & 1 & 1179648 & 44.05 & 131 & Sample replay \\
64 & 1.778 & 2097152 & 72.05 & 131 & Sample replay \\
72 & 2 & 2359296 & 80.05 & 131 & Sample replay \\
128 & 3.556 & 4194304 & 136.1 & 131 & Sample replay \\
256 & 7.111 & 8388608 & 264.1 & 131 & Sample replay \\
512 & 14.22 & 16777216 & 520.1 & 131 & Sample replay \\
1024 & 28.44 & 33554432 & 1032 & 132 & Sample replay \\
2048 & 56.89 & 67108864 & 2056 & 136 & Sample replay \\
4096 & 113.8 & 134217728 & 4104 & 144 & Sample replay \\
8192 & 227.6 & 268435456 & 8200 & 160 & Sample replay \\
9978 & 277.2 & 326959104 & 9986 & 169 & Sample replay \\
\end{longtable}
\endgroup

All epochs compare complete aggregate digests across access modes, traversal patterns and trials. Full CPU replay is restricted to the small sets flagged above; larger sets use sampled independent word replay. Digest agreement can collide and is not a byte-for-byte readback proof of every large buffer. Exact seed/state byte counts, free-memory observations and source receipts are retained in the summary JSON.

\subsection{Every word-cache timing trial}

Each cell lists trials in increasing trial order. Times are milliseconds for 3 measured epochs per trial. GPU sums exclude reset/warmup, digest readback and sampled verification. Host wall includes measured epochs, digest aggregation/readback and per-epoch sampled verification; reset/warmup are outside the timer. Raw per-epoch values remain in the JSON.

\Needspace{6\baselineskip}

\paragraph{Streaming}

\begingroup
\scriptsize
\setlength{\tabcolsep}{2.5pt}
\renewcommand{\arraystretch}{1.16}
\begin{longtable}{rrp{2.85cm}p{2.85cm}p{2.85cm}p{2.85cm}}
\toprule
MiB & L2 ratio & Texture GPU & Global GPU & Texture wall & Global wall \\
\midrule\endfirsthead
\toprule
MiB & L2 ratio & Texture GPU & Global GPU & Texture wall & Global wall \\
\midrule\endhead
\bottomrule\endfoot
0.0625 & 0.001736 & \shortstack[r]{0.0943\\0.06342\\0.1115} & \shortstack[r]{0.09088\\0.08512\\0.08682} & \shortstack[r]{0.4775\\0.3175\\0.4189} & \shortstack[r]{0.4435\\0.3453\\0.4072} \\
0.25 & 0.006944 & \shortstack[r]{0.103\\0.1041\\0.08954} & \shortstack[r]{0.08291\\0.09619\\0.06877} & \shortstack[r]{0.411\\0.4149\\0.4951} & \shortstack[r]{0.4639\\0.4567\\0.3928} \\
1 & 0.02778 & \shortstack[r]{0.1001\\0.08781\\0.06947} & \shortstack[r]{0.08506\\0.1161\\0.0887} & \shortstack[r]{0.4529\\0.4355\\0.3398} & \shortstack[r]{0.4391\\0.3581\\0.4339} \\
4 & 0.1111 & \shortstack[r]{0.1098\\0.1087\\0.0975} & \shortstack[r]{0.09859\\0.09677\\0.1707} & \shortstack[r]{0.4743\\0.4717\\0.3979} & \shortstack[r]{0.4767\\0.4536\\0.5152} \\
16 & 0.4444 & \shortstack[r]{0.2865\\0.2964\\0.2933} & \shortstack[r]{0.2965\\0.2728\\0.2816} & \shortstack[r]{0.8373\\0.8681\\0.8343} & \shortstack[r]{0.8749\\0.8475\\0.821} \\
18 & 0.5 & \shortstack[r]{0.4168\\0.3638\\0.3957} & \shortstack[r]{0.4044\\0.3873\\0.3713} & \shortstack[r]{1.133\\0.9255\\1.086} & \shortstack[r]{1.075\\1.069\\1.033} \\
36 & 1 & \shortstack[r]{0.6731\\0.6868\\0.7426} & \shortstack[r]{0.6843\\0.758\\0.692} & \shortstack[r]{1.612\\1.793\\1.987} & \shortstack[r]{1.856\\2.011\\1.857} \\
64 & 1.778 & \shortstack[r]{1.334\\1.242\\1.243} & \shortstack[r]{1.256\\1.149\\1.136} & \shortstack[r]{3.259\\2.858\\2.868} & \shortstack[r]{3.063\\2.814\\2.79} \\
72 & 2 & \shortstack[r]{1.368\\1.428\\1.4} & \shortstack[r]{1.286\\1.362\\1.35} & \shortstack[r]{3.295\\3.367\\3.236} & \shortstack[r]{3.029\\3.215\\3.174} \\
128 & 3.556 & \shortstack[r]{2.508\\2.599\\2.534} & \shortstack[r]{2.394\\2.435\\2.349} & \shortstack[r]{5.562\\5.689\\5.66} & \shortstack[r]{5.457\\5.465\\5.613} \\
256 & 7.111 & \shortstack[r]{4.955\\5.044\\5.356} & \shortstack[r]{4.758\\4.807\\5.128} & \shortstack[r]{10.89\\10.91\\11.54} & \shortstack[r]{10.89\\10.73\\11.7} \\
512 & 14.22 & \shortstack[r]{10.62\\9.883\\9.796} & \shortstack[r]{10.02\\9.388\\10.79} & \shortstack[r]{23.77\\21.51\\21.02} & \shortstack[r]{22.34\\21.1\\22.74} \\
1024 & 28.44 & \shortstack[r]{19.93\\19.58\\20.1} & \shortstack[r]{18.97\\20.1\\20.27} & \shortstack[r]{43.12\\42.99\\43.45} & \shortstack[r]{42.31\\45.29\\45.11} \\
2048 & 56.89 & \shortstack[r]{40.27\\39.72\\40.54} & \shortstack[r]{37.88\\39.69\\37.34} & \shortstack[r]{88.64\\86.52\\88.7} & \shortstack[r]{84.86\\88.85\\83.51} \\
4096 & 113.8 & \shortstack[r]{78.44\\79.82\\80.53} & \shortstack[r]{75.78\\75.06\\74.62} & \shortstack[r]{170.9\\173.5\\173.7} & \shortstack[r]{167.6\\166.9\\167.5} \\
8192 & 227.6 & \shortstack[r]{176.1\\176.6\\176.4} & \shortstack[r]{168.7\\168.5\\170.9} & \shortstack[r]{360.7\\361.6\\361} & \shortstack[r]{353.9\\354.6\\356.2} \\
9978 & 277.2 & \shortstack[r]{215.2\\217.7\\217.2} & \shortstack[r]{205.6\\205.6\\203.9} & \shortstack[r]{441.7\\443.3\\445.3} & \shortstack[r]{432.3\\432.8\\428.2} \\
\end{longtable}
\endgroup

\Needspace{6\baselineskip}

\paragraph{Seeded coprime affine permutation}

\begingroup
\scriptsize
\setlength{\tabcolsep}{2.5pt}
\renewcommand{\arraystretch}{1.16}
\begin{longtable}{rrp{2.85cm}p{2.85cm}p{2.85cm}p{2.85cm}}
\toprule
MiB & L2 ratio & Texture GPU & Global GPU & Texture wall & Global wall \\
\midrule\endfirsthead
\toprule
MiB & L2 ratio & Texture GPU & Global GPU & Texture wall & Global wall \\
\midrule\endhead
\bottomrule\endfoot
0.0625 & 0.001736 & \shortstack[r]{0.08128\\0.08371\\0.07434} & \shortstack[r]{0.08086\\0.09946\\0.06886} & \shortstack[r]{0.3697\\0.3995\\0.2667} & \shortstack[r]{0.4077\\0.4421\\0.3064} \\
0.25 & 0.006944 & \shortstack[r]{0.08173\\0.05789\\0.06675} & \shortstack[r]{0.1023\\0.08483\\0.07347} & \shortstack[r]{0.4128\\0.3696\\0.2881} & \shortstack[r]{0.3091\\0.439\\0.2909} \\
1 & 0.02778 & \shortstack[r]{0.08902\\0.08368\\0.08973} & \shortstack[r]{0.08717\\0.08528\\0.1067} & \shortstack[r]{0.3428\\0.3707\\0.4084} & \shortstack[r]{0.4495\\0.3998\\0.4131} \\
4 & 0.1111 & \shortstack[r]{0.1156\\0.132\\0.1162} & \shortstack[r]{0.1452\\0.1134\\0.09738} & \shortstack[r]{0.4454\\0.4357\\0.4114} & \shortstack[r]{0.4714\\0.4211\\0.3176} \\
16 & 0.4444 & \shortstack[r]{0.3593\\0.3602\\0.3588} & \shortstack[r]{0.3567\\0.3686\\0.3426} & \shortstack[r]{0.8998\\0.9185\\0.9292} & \shortstack[r]{0.908\\0.92\\0.9077} \\
18 & 0.5 & \shortstack[r]{0.4314\\0.4907\\0.4392} & \shortstack[r]{0.4838\\0.4725\\0.426} & \shortstack[r]{1.153\\1.192\\1.15} & \shortstack[r]{1.149\\1.122\\1.072} \\
36 & 1 & \shortstack[r]{1.045\\1.063\\1.027} & \shortstack[r]{1.054\\1.089\\1.092} & \shortstack[r]{2.173\\2.233\\2.189} & \shortstack[r]{2.079\\2.311\\2.311} \\
64 & 1.778 & \shortstack[r]{5.074\\5.005\\5.017} & \shortstack[r]{4.98\\4.96\\4.928} & \shortstack[r]{6.78\\6.69\\6.639} & \shortstack[r]{6.718\\6.592\\6.621} \\
72 & 2 & \shortstack[r]{5.314\\5.344\\5.374} & \shortstack[r]{5.328\\5.317\\5.306} & \shortstack[r]{7.23\\7.39\\7.244} & \shortstack[r]{7.21\\7.174\\7.139} \\
128 & 3.556 & \shortstack[r]{11.15\\11.09\\11.13} & \shortstack[r]{11.04\\10.93\\11} & \shortstack[r]{14.32\\14.28\\14.52} & \shortstack[r]{14.22\\14.09\\14.09} \\
256 & 7.111 & \shortstack[r]{23.14\\23.16\\23.06} & \shortstack[r]{22.84\\22.78\\23.13} & \shortstack[r]{29.41\\29.29\\28.97} & \shortstack[r]{28.76\\28.82\\29.33} \\
512 & 14.22 & \shortstack[r]{47.34\\46.82\\46.84} & \shortstack[r]{46.38\\46.47\\46.74} & \shortstack[r]{59.75\\58.81\\58.99} & \shortstack[r]{58.67\\58.28\\59.4} \\
1024 & 28.44 & \shortstack[r]{96.93\\97.79\\97.4} & \shortstack[r]{96.6\\96.18\\97.04} & \shortstack[r]{121.7\\122.6\\121.5} & \shortstack[r]{121.2\\120.4\\120.5} \\
2048 & 56.89 & \shortstack[r]{195.9\\194.2\\196.1} & \shortstack[r]{192.7\\193.5\\193} & \shortstack[r]{244\\241.9\\243.1} & \shortstack[r]{240.4\\241.5\\240.3} \\
4096 & 113.8 & \shortstack[r]{387.6\\388.3\\387.3} & \shortstack[r]{386.7\\386.5\\387.4} & \shortstack[r]{481.8\\481.8\\481.5} & \shortstack[r]{481\\480.7\\484.5} \\
8192 & 227.6 & \shortstack[r]{798\\776.6\\797.9} & \shortstack[r]{774.2\\799\\774.3} & \shortstack[r]{985.8\\966.8\\988.9} & \shortstack[r]{962.9\\988\\965.8} \\
9978 & 277.2 & \shortstack[r]{980.1\\965.2\\984} & \shortstack[r]{942.7\\959.1\\943.8} & \shortstack[r]{1211\\1198\\1219} & \shortstack[r]{1174\\1193\\1177} \\
\end{longtable}
\endgroup

Texture beat global loads in 7 of 34 GPU-median comparisons and 9 of 34 wall-median comparisons. Cache hit rates were not measured and no persisting-L2 window was set. Logical GB/s is workload throughput, not measured physical DRAM bandwidth or a universal texture advantage.

\paragraph{Reproducible source receipts.}

Measured benchmark source commit: \path{4d8992292a7fa0fc43fc050ec090512552f7d067}. Subsequent source edits are not implied to have been measured by this record.

All trials, setup timings and metadata are preserved in \path{review/resident_summary.json}. Its inputs are \path{review/resident_full.json} and \path{review/word_cache_full.json}.

\begin{small}\noindent Resident SHA-256: \path{c8c25de5a13b9f559c65081ebd597547c1a5d41753edc09deff7db858fd4a94d}.\par

\noindent Word-cache SHA-256: \path{adb44ca24fbab2e4534bae716ca94213bd28ba4cf1908ae4d62fc0b0fbf93946}.\end{small}
}{%
\begin{panel}
The resident and working-set run record is pending. No timing, cache-hit
ratio, completed working-set size or new correctness count is asserted by
this mathematical scaffold.
\end{panel}}
